\chapter{Fast multipole field evaluation}\label{chap:fmm}

\section{Purpose and scope}
The fast multipole method (FMM) accelerates the application of a nonlocal field
operator without assembling its dense interaction matrix. For source strengths
$\bm q_s$ and target positions $\bm x_t$, the basic operation is
\begin{equation}
 \bm u_t=\sum_s \mathsf K(\bm x_t-\bm y_s)\bm q_s.
 \label{eq:fmmDiscreteOperator}
\end{equation}
The kernel $\mathsf K$ may be scalar or tensor valued. Dislocation stress and
mobile-species concentration use different kernels and source strengths, but
share the geometry, hierarchy and translation machinery. In a climb solve,
each GMRES iteration changes the strengths; it does not require reconstruction
of the geometry or an explicit dense Green's-function matrix.

The implementation uses tensor-product interpolation coefficients as its
multipole representation. These coefficients encode the effect of all sources
in a box on an interpolation basis. They play the role of moments, but are not
a spherical-harmonic or Cartesian Taylor expansion of a particular kernel.
This makes it possible to register different kernel families with the same
FMM framework.
\sourcefile{include/FMM/FMMManager.h}
\sourcefile{include/FMM/KokkosParticleFMM.h}

\section{A lattice-aligned box hierarchy}
Write a physical point as
\begin{equation}
 \bm x=\bm x_0+\mathsf H\bm\xi,\qquad
 \mathsf H=[\bm a_1\ \bm a_2\ \bm a_3],\qquad
 \bm\xi\in[0,1]^3.
 \label{eq:fmmAffineCell}
\end{equation}
The columns of $\mathsf H$ are the physical root-cell edges. The hierarchy is
axis aligned in fractional coordinates $\bm\xi$; its physical cells may be
skewed parallelepipeds. At level $d$, each coordinate interval has been divided
into $2^d$ intervals, giving $8^d$ boxes with edge vectors
$2^{-d}\bm a_i$. Every parent has eight children. The current hierarchy is
uniform, rather than adaptively refined by particle occupancy.

\code{fmmTreeDepth} selects the leaf level. Increasing it reduces the number
of sources per leaf and the direct work per target, while increasing the number
of boxes, interpolation coefficients and translations. Empty boxes need not
contribute source work, but the uniform hierarchy still has a memory cost.

Without periodic directions, the root retains the Cartesian mesh bounds.
With periodic directions, the corresponding cell edges follow the mesh's
periodic lattice basis. For partial periodicity, an orthogonal complement
supplies the remaining directions. Bounds in this basis are obtained from the
shared grain polyhedron vertices, not from every finite-element mesh node.
Each fractional coordinate is an affine function of physical position, so its
extrema over a convex grain occur at vertices. Taking extrema over all grain
vertices therefore encloses the entire polycrystal. Shared vertices are read
once through the polyhedral vertex factory.

This construction provides a covering cell aligned with the periodic lattice.
It must not be confused with an automatic identification of opposite tree
faces or with a periodic-image summation. In particular, the covering cell of
a general partially periodic domain need not itself be a fundamental periodic
cell.
\sourcefile{include/FMM/BoxTree.h}
\sourcefile{include/FMM/LatticeBoxTree.h}
\sourcefile{src/DefectiveCrystal/DefectiveCrystalBase.cpp}

\subsection{Source and target assignment}
Source and field roles are distinct. A segment can contribute a field even
when no velocity or force is required on that segment. Sources and targets
are assigned to leaves using fractional coordinates, while kernel evaluation,
Burgers vectors, stresses, diffusion tensors and physical line lengths remain
in Cartesian coordinates. An affine coordinate map for the tree does not
change the physical kernel.

For segment-based evaluation, straight segments are split at leaf faces.
A segment crossing several boxes contributes one piece to each intersected
box. Splitting changes the integration bookkeeping, not the physical network.
The shared split cache is rebuilt when segment geometry or the tree changes;
stress, concentration and climb reuse the same pieces. Roundoff-equivalent
crossings at a box edge or vertex are merged to avoid spurious tiny pieces.

Source coefficients are integrated along each piece. The interpolation basis
restricted to a straight line is polynomial, so suitable Gaussian rules
integrate these coefficients exactly up to floating-point error. This does
not assert that Gaussian quadrature integrates the physical near-field kernel
exactly. Near interactions use the corresponding segment kernel, including
endpoint/table formulas where available.
\sourcefile{include/DislocationDynamics/SegmentSplitCache.h}
\sourcefile{include/DislocationDynamics/StraightDislocationSegmentMoments.h}

\section{Near interactions and physical admissibility}
Index-space proximity is insufficient for a skewed cell: two boxes separated
by several indices can still be close in physical space. The implementation
uses a conservative enclosing-sphere criterion. For congruent level-$d$ boxes,
let
\begin{equation}
 D_d=2^{-d}\max_{\bm\epsilon\in\{-1,1\}^3}
          \|\mathsf H\bm\epsilon\|.
\end{equation}
This is the maximum physical diagonal and twice the enclosing-sphere radius.
Boxes with centers $\bm c_t,\bm c_s$ are accepted for far translation when
\begin{equation}
 \|\bm c_t-\bm c_s\|>1.01D_d.
 \label{eq:fmmAdmissibility}
\end{equation}
Otherwise the traversal descends to their children. At the leaf level the
unaccepted pairs form the near list. The same criterion determines the direct
segment neighborhood, avoiding gaps or overlap between near and far work.

For cubic cells this gives the familiar self box and 26 neighbors. Skewed or
elongated cells can require more neighbors. A strongly distorted lattice is
therefore supported conservatively but may increase direct work considerably.
The geometric acceptance test is not itself a prescribed field-error
 tolerance; interpolation order and kernel smoothness still determine accuracy.
\sourcefile{include/FMM/LeafBoxNeighbors.h}

\section{How the hierarchy evaluates the far field}
For interpolation order $p$, each box has $p^3$ interpolation nodes. The
calculation consists of the following passes:
\begin{enumerate}
 \item \textbf{Source-to-multipole (P2M).} Construct leaf coefficients from
 particles or integrate them directly over segment pieces.
 \item \textbf{Multipole-to-multipole (M2M).} Transfer child coefficients to
 parents. A parent represents all sources in its descendants.
 \item \textbf{Multipole-to-local (M2L).} Translate accepted source-box
 coefficients to a target-box local field. A dual-tree traversal accepts a
 pair once or subdivides it; distant groups are handled at coarse levels.
 \item \textbf{Local-to-local (L2L).} Propagate accumulated local fields from
 parents to children.
 \item \textbf{Local-to-target (L2P).} Interpolate the leaf local field at
 each target and add the required near contribution.
\end{enumerate}
Thus distant particles do not interact individually with every target.
Their aggregate coefficients are reused across many targets and across the
hierarchy. The source representation is independent of the current target
values; in a linear solve only strengths and the resulting coefficients need
updating while geometry is fixed.

For a translation-invariant kernel, boxes at the same level with the same
relative lattice offset share an M2L operator. Cached matrices depend on the
cell geometry, interpolation order and kernel parameters, not on individual
particle locations. Grouping multiple source-box coefficient vectors forms
matrix--matrix products. The current CPU path uses Eigen for these grouped
products; the surrounding FMM passes use Kokkos. A grouped GPU implementation
is deferred. Cache budgets can limit which translation matrices are retained.

\section{Unsplit and split kernel families}
\subsection{Ordinary FMM}
With \code{multipoleMode=FMM}, the full core-regularized kernel is used in both
near and far interactions:
\begin{equation}
 \bm u_t\simeq
 \sum_{s\in\mathcal N_t}\mathsf K(\bm x_t-\bm y_s)\bm q_s
 +\operatorname{FMM}_{\mathcal F_t}[\mathsf K\bm q].
\end{equation}
Here $\mathcal N_t$ and $\mathcal F_t$ partition the source interactions.
The chosen \code{coreType}, currently Cai or Helmholtz, specifies the physical
regularization. Ordinary FMM does not introduce an additional screening length.
With \code{multipoleMode=none}, the accelerated evaluation is bypassed and the
selected core's direct segment evaluation is used instead.

\subsection{Ewald-split FMM}
With \code{multipoleMode=ewaldFMM}, the framework uses a decomposition
\begin{equation}
 \mathsf K=\mathsf K_{\mathrm{near}}+\mathsf K_{\mathrm{far}},
 \qquad
 \bm u_t\simeq
 \sum_{s\in\mathcal N_t}\mathsf K_{\mathrm{near}}(\bm x_t-\bm y_s)\bm q_s
 +\operatorname{FMM}_{\mathrm{all}}[\mathsf K_{\mathrm{far}}\bm q].
 \label{eq:fmmSplit}
\end{equation}
The near kernel decays rapidly and is omitted outside the near-box stencil.
The smoother far kernel contributes from \emph{all} source boxes, including
the self box and near neighbors. Those additional smooth translations are
essential: excluding them would omit part of the full interaction. Cached
self/neighbor translations avoid evaluating the smooth kernel separately for
every nearby source--target pair.

The current Ewald path uses the Helmholtz core and screening formulation.
The physical core length $\ell$ and numerical Ewald width $h$ have different
roles. The former regularizes the physical field; the latter controls the
split and should not alter the full field in an exact evaluation. The input
\code{ewaldWidthFactor} determines $h$ from the minimum leaf face spacing:
\begin{equation}
 h=\texttt{ewaldWidthFactor}\;2^{-d}
       \min_i\frac{1}{\|\text{row}_i(\mathsf H^{-1})\|}.
\end{equation}
A smaller width makes the short field more local but leaves a less smooth far
field to interpolate. Larger widths improve smoothness while increasing the
short tail omitted beyond the stencil. Short-field truncation and long-field
interpolation are separate errors and must be measured separately. Ewald FMM
optimization is currently deferred; ordinary FMM with the Cai core is the
working default.

\subsection{A common interface}
Kernel families expose full, near and far evaluations and declare whether
they are split. For an unsplit family, near and far dispatch to the full
kernel. For a split family, they provide the separate contributions; the full
kernel retains its direct expression rather than being evaluated by splitting
and recombining it. Segment-specific stress and concentration formulas live
in the corresponding Cai, Helmholtz and Helmholtz--Ewald classes, with
polynomial moment integration in a shared helper. New physics should supply
its kernel and source representation while reusing the hierarchy and plans.
\sourcefile{include/DislocationDynamics/StraightDislocationSegmentBase.h}
\sourcefile{include/DislocationDynamics/SegmentConcentrationPlan.h}

\section{Periodicity: geometry versus image summation}
For $r$ active periodic vectors collected in $\mathsf A$, the intended periodic
field has the form
\begin{equation}
 \bm u(\bm x)=\sum_{\bm n\in\mathbb Z^r}\sum_s
 \mathsf K(\bm x-\bm y_s-\mathsf A\bm n)\bm q_s,
\end{equation}
subject to the kernel's compatibility conditions and summation convention.
For partial periodicity, translations occur only in the active directions;
physical boundaries in the remaining directions are not wrapped.

The lattice-aligned tree implements the geometry needed for this extension.
It does \emph{not} yet implement the periodic-image field. Periodic FMM
requests remain explicitly rejected rather than silently returning free-space
results. The next layer must account for nearby translated boxes and the
remaining image lattice, without double counting images or central-cell
interactions. A finite image sum must be validated against the same finite
image set in direct evaluation. It must not be presented as a converged
infinite periodic sum without a convergence study. Infinite sums additionally
require the appropriate neutrality/compatibility conditions and macroscopic
or zero-mode convention for each kernel.

This separation also matters for FEM coupling: periodic identification of
mesh degrees of freedom, periodic continuation of the dislocation network,
and periodic evaluation of its fields are distinct operations. Aligning the
FMM tree solves their geometric coordination, not the field boundary-value
problem by itself.
